% treatment.tex — Full treatment

\section{Treatment}

\subsection{Movement I: The Father}

The film opens with an image of absolute precision: a thin man on a Berlin street in winter, coughing blood outside a butcher's window. The blood on his hand. The meat in the window. A woman passing who does not notice. This is Franz Kafka in 1923, and this image --- the body's refusal, the world's indifference, the clinical proximity of human and animal flesh --- will resonate through everything that follows.

We do not stay in Berlin. We slide backward, almost without permission, to Prague in 1912. The past arrives not as flashback but as contamination: a sound, a cut, and we are in the Kafka family apartment on the Altstädter Ring, where the air is thick with coffee, expectation, and the particular gravity of Hermann Kafka.

Hermann is the film's psychic center of gravity even when he is not on screen. He is not a monster. He is something worse: a comprehensible man. A self-made businessman, practical and loud and alive with appetite. He fills rooms. He fills his son's life. His authority is not tyrannical but atmospheric --- the pressure of normality, of expectation, of a father who cannot understand why his brilliant son sits up all night writing instead of sleeping like a stone and advancing at the office.

We meet Julie, the mother who mediates and protects no one. Ottla, the youngest sister, Kafka's only real ally. We see the family apartment as Kafka experiences it: heavy, furnished, suffocating in its respectability.

The centerpiece of this movement is the shop scene. Hermann in his fancy goods store --- ribbons, buttons, gloves --- is magnificent. He knows every customer, every margin. When Franz arrives, Hermann's attention shifts: the same assessment he gives his merchandise, applied to his son. The casual humiliation is not dramatic. It is atmospheric. A posture correction. A remark about eating. The introduction to a customer that manages to convey, without ever stating it, that Franz is a disappointment in his father's specific, practical terms.

The movement ends in the family apartment at night. Kafka at his desk, writing --- the only place where the architecture of his inner life can take shape. Then the corridor. He walks to the bathroom through the dark apartment. The hallway keeps going. His father breathes behind one of the doors. Then the corridor is normal again. But in the mirror, behind him, it still isn't.

\subsection{Movement II: The Office}

The Workers' Accident Insurance Institute is the film's second great arena. Kafka's day job was not incidental to his art --- it was its dark mirror. Here, human suffering is processed into paper. Crushed hands become claim forms. Factory deaths become statistical categories. Kafka was genuinely good at this work, which makes it worse.

The signature scene of this movement is the elevator. It is a small scene --- perhaps three minutes on screen --- but it must play like a metaphysical event. Kafka waits for the elevator with three colleagues. The elevator arrives. The three men enter. There is room for one more. Kafka steps forward. The operator closes the gate. The elevator rises. Kafka is left standing alone. Nothing has happened. Everything has happened.

We see Kafka at his desk, reading the Kraus sawmill file --- a worker who lost three fingers because a safety guard was removed to save four seconds per board. Kafka inspects the factory personally. His competence is striking and revealing: a man who understands systems better than anyone, and who knows that understanding a system does not free you from it. The Kraus case will recur across movements, a persistent proof that the world's machinery does not learn.

When offered a promotion to senior assessor, Kafka declines. The additional hours would cost him his nights --- the only hours that belong to him. His director is baffled. Kafka is, for once, clear about what he wants: not advancement, but time.

Max Brod enters. In a Prague café, over coffee and manuscripts, we see the friendship that will define Kafka's legacy. Brod is everything Kafka isn't: sociable, optimistic, certain of his friend's genius. Their dynamic is affectionate and slightly abrasive --- Brod pushes, Kafka retreats, and beneath the affection, a philosophical disagreement is already forming. Brod believes literature belongs to the world. Kafka believes it belongs to the writer. This fault line will crack open in the final act.

The movement ends with a corridor. Kafka walking through the Institute after hours. The corridor, which should have twelve doors, has twenty-three. The door numbers morph into claim file numbers. Reality bends, briefly, almost imperceptibly. Then snaps back. Kafka adjusts his hat and goes home.

\subsection{Movement III: The Women}

This movement belongs to the two women who most precisely reveal Kafka's terror of intimacy --- Felice Bauer and Milena Jesenská. Each represents a different version of the life he cannot enter, and each exposes a different failure in him.

Felice comes first. She represents the one thing Kafka most desires and most dreads: an ordinary life. She is practical, organized, German-Jewish, a working woman at a dictation machine company in Berlin. Critically, we see her at work --- brisk, competent, entirely at ease inside a modern office world of machines, orders, and timetables. She is not an abstraction called normal life; she is the actual adult life Kafka is trying and failing to enter. Their courtship is conducted largely by letter --- Kafka writes to her obsessively, in a register of emotional exposure he would never permit in person.

When they finally reach the point of engagement, the dinner scene --- both families present, Hermann booming about appetite and structure, expectations crystallized into social performance --- plays like a tribunal. Kafka sits at the center of the table and smiles. The smile is the greatest performance of his life. Hermann claps him on the shoulder: ``My boy has finally found some appetite!'' Something behind the smile collapses.

But the crucial revision is what follows. Felice does not simply receive Kafka's withdrawal. She acts. She has read his letters to Grete Bloch --- letters in which he describes, with devastating eloquence, the impossibility of the life she is offering. She confronts him in a hotel room, arriving from her office with a return train ticket already in hand. She ends it. ``You prefer wanting a thing to having it.'' She leaves. Kafka sits alone. For once, the desk offers no escape.

The aftermath: Kafka returns to the office. A new development in the Kraus sawmill case --- a second worker injured on the same machine. The world's machinery continues regardless of his emotional state. He picks up his pen. The best writing he will do that day is an insurance report.

Then Milena --- six years later, in the same movement, because the architecture of this film is emotional, not calendrical. Where Felice was social adulthood --- schedules, office light, ordinary life --- Milena Jesenská is intellectual and erotic recognition at a level Kafka has never survived. She translates his work into Czech and, through the translations, reads him as no one else can --- not just his fragility but his evasions, his arrangements, the structure of his retreat. Their meeting in Vienna is electric and impossible. ``Because with you I thought I might fail in person and remain anyway,'' he tells her. She accepts this. She does not accept it happily. The movement ends on a Vienna train platform, Milena refusing to wave, having understood everything about him except whether he actually means to live.

\subsection{Movement IV: The Illness}

In 1917, Kafka coughs blood. The echo of the film's opening image is precise and devastating. A doctor confirms tuberculosis. Ottla takes him to Zürau, where the country air and the absence of Prague produce the sharpest prose of his life --- aphorisms like stones. Hermann arrives unannounced with a basket from Julie --- cheese, bread, smoked meat. For one scene, a father and a son eat bread in a small room. It is almost ordinary. It is almost enough.

But the office does not stay in Prague. Marschner sends the Kraus appeal file to Zürau with a note asking for one more memorandum. Kafka answers from sick leave. The moment matters because it preserves one of the film's most distinctive truths: even in illness, he remains one of the few men in his world capable of forcing an institution to say exactly what happened and who paid for the savings in blood. Then he sets that memorandum aside and turns to the unwinnable case. The office beat now serves the Letter by contrast: exact prose for strangers, then impossible prose for the father.

The centerpiece of this movement --- and of the film --- is the Letter to His Father. Written in 1919, over a hundred pages, it is simultaneously a confession, an accusation, a legal brief, and a love letter. Kafka stands at his father's study door, the manuscript in his hands, his hand on the handle, his father's breathing on the other side. He cannot open it. He walks to the kitchen instead and gives the manuscript to his mother with a request to deliver it. But Julie refuses --- not through cowardice, but through the painful clarity of a woman who has spent thirty years managing the space between two people who cannot survive each other's truth. ``If I carry this into that room, he will hear accusation before he hears pain,'' she tells Franz. She puts it in a drawer. Hermann Kafka will never read it.

The letter is presented not as continuous voiceover but as fragments --- heard over images of Kafka's face, his father's hands, the apartment, the shop. The emotional force builds through accumulation.

Ottla confronts Franz about the letter with a directness no one else in his life possesses: ``You wrote a hundred pages because paper doesn't interrupt. Then you gave them to Mother because she does.'' The truth of this assessment is devastating.

Brod visits the sanatorium. He is worried about the manuscripts, and the philosophical disagreement that was seeded in their earlier conversations now becomes explicit. Brod believes the work belongs to the world; Kafka believes it belongs to the writer. The conversation about destruction begins here, gently. It will end as a command.

\subsection{Movement V: The Fire}

Dora Diamant enters the film like a window opening in a sealed room. She is 25, warm, practical, from a Hasidic family in Poland. Kafka meets her at a summer camp on the Baltic coast. Where Felice offered structure and Milena offered dangerous recognition, Dora offers presence: the body fed, the room shared, the day actually lived. She is the first person in his life with whom he is simply, uncomplicatedly present.

They move to Berlin. The apartment is cold, small, poor. They are happy. The happiness is real and brief and almost unbearable to watch, because we know what is coming. For the first time, Kafka hears himself think. ``And what do you think?'' Dora asks. ``That I'm happy,'' he answers, diagnosing the condition like a man who has read about it but never personally encountered it.

In Berlin, Julie mails back the undelivered Letter to His Father. The packet enters the apartment like a curse that has finally located its proper address. Kafka does not reread it. He puts it in a drawer that will not shut cleanly. The father's pressure has followed him even here, into the only room that ever felt like his own.

Marschner writes too. The Kraus appeal has failed. Compensation stands. Guards must be installed across all three mills. It is not triumph exactly --- too many hands have already paid for the lesson --- but it proves that Kafka's office life was not merely metaphor. His sentences altered material conditions in the world, even if only late, even if never enough.

Kafka burns manuscripts with Dora. The fire is both destruction and completion. But he does not burn the Letter. It is not unfinished; it is undeliverable. That distinction matters to him. Then Brod visits, and the moral center of the film arrives: the Burn Everything scene. What distinguishes this version is the depth of philosophical disagreement. Brod argues that destruction is its own vanity --- the writer who insists his unfinished work be destroyed is still curating, still controlling. Kafka counters that prayers are private, that a half-built house belongs to the person who couldn't finish building it. Neither is wrong. Both know the other will not yield. Brod agrees. He will not do it.

Meanwhile, Kafka writes to Dora's Hasidic father requesting permission to marry --- and the echo of the Letter to His Father is explicit. He has written to a father before. A hundred pages, never read. This letter is shorter, asks for less. It is also refused. Another form returned, stamped rejected.

Kafka's health collapses. He can no longer stay in Berlin.

\subsection{Movement VI: The Last Room}

The sanatorium at Kierling, near Vienna. Kafka's tuberculosis has reached his larynx. He cannot speak. He cannot eat. He communicates by handwritten notes on slips of paper.

The notes are funny, precise, heartbreaking. He corrects the proofs of \textit{A Hunger Artist} --- his last published collection, whose title story concerns a man who starves because he never found food he liked. The parallel with Kafka's own condition is total and unstated. Even here, the father remains in the frame: Kafka begins to write instructions to send a copy, stops, and leaves the line crossed out. The hundred-page letter could never cross the room it was written for. Perhaps this small book could.

The final scene is the emotional summit of the film. Kafka lies in bed, physically devastated, eyes alive. He writes notes to Klopstock: ``Kill me, or else you are a murderer.'' He writes to Hoffmann: ``I have been patient my entire life. Do not ask me to be patient with this.'' Then, to all of them: ``No fear now.''

He dies on June 3, 1924, with Dora and Klopstock present.

Coda: Julie waits in the Prague apartment with two things for Brod --- a box of Franz's notebooks and the thick sealed packet addressed ``For Father. Via Mother.'' She has kept it since he put it in her hands. She has kept it from the drawer Hermann nearly opened. She tells Brod to do what his conscience can survive with the notebooks, but the letter stays tied --- if Franz had wanted it opened, he knew Brod's address. Brod takes them both.

At his own desk, surrounded by the manuscripts, Brod picks up a matchbox and sets it beside \textit{The Trial}. He reads the first page. Then he puts the matchbox in a drawer and pulls the next manuscript toward him.

Final image: the Berlin apartment, empty now. Moonlight on the small table. A Hebrew primer, one chipped bowl, the folded registration paper with its purple stamp, and a small brass key. The drawer is shut. For the first time, it closes all the way. Hold.
